\documentclass[12pt,a4paper]{article}
\usepackage{amsfonts, amssymb, amsmath}
\usepackage{fullpage}
\usepackage{parskip} % skip a line instead of indenting
% \usepackage{graphicx} % for inserting images
\usepackage{amsthm}
\usepackage{xcolor}
\usepackage{tikz} % for plot

\newcommand{\remark}[1]{
    $>$ {\color{blue} #1}
}

\begin{document}

\section*{Gossip Principle Analysis}

Imagine the network as a population of $n+1$ individuals:
\begin{itemize}
    \item $y$ (Infected): Nodes that have received the information.
    \item $x$ (Uninfected): Nodes that have not yet received the information
    \item $x + y = n + 1$: The total population remains constant.
    \item $\beta$ (Contact Rate): The probabillity that two nodes meet and the message is transmitted.
\end{itemize}

Init $t = 0$: Only 1 node knows the news ($y_0$ = 1), and $n$ nodes are unaware ($x_0$ = $n$).

We wan to know

\[\frac{dx}{dt} = -\beta xy\]

Why $-\beta xy$?

First, any pair $(x, y)$ that meets will result in a message transmission. 
There are totally $xy$ possible pairs and each pair has a $\beta$ chance to transmit the message; 
the negative sign indicates that $x$ should decrease over time.

\[\Rightarrow y = \frac{n+1}{1 + n \cdot e^{-\beta(n+1)t}} \quad \text{(S curve)}\]

However, In a Gossip protocal, we use ``active selection'' rather than ``random collision'':

Each node randomly pick $b$ targets every round (aka. Fan-out). Since there are $n$ potential targets, the contact rate is $\beta \approx \frac{b}{n}$. 
According to the most powerful feature of the Gossip protocol, the time required for informantion to disseminate across the entire network is $O(\log n)$, 
we set time as $t = c\log n$ (where $c$ is a constant).

\[\Rightarrow y \approx \frac{n+1}{1 + n \cdot n^{-bc}} = \frac{n+1}{1 + n^{1-bc}}\]


Key Takeaways:

\begin{enumerate}
    \item Efficiency: Reaches the entire network in $O(\log n)$ rounds.
    \item Scalability: Each node carries a constant load of $b$ messages per round, regardless of network size.
\end{enumerate}

\end{document}